\begin{abstract}

An octree is a structure which partitions a 3D space into octants recursively based on positions of objects within the domain. This structure is heavily utilized in scientific simulation, computer graphics, and robotics. One major use case of an octree is in the field of particle simulations. Simulations calculate the particle interactions between millions, billions, or even trillions of particles, and this complexity of calculating per-pair interactions becomes prohibitive for the system.

Recent works have expanded upon traditional octree construction methods to allow for distributed simulation on graphics processing units (GPUs). We explore and benchmark a novel method called Cornerstone ~\cite{cornerstone} that implements these methods and develops a construction method that allows an octree to be rebalanced on-the-fly which is favorable for particle simulations or any problem which frequently perturbs the positions of the points in some way. This optimization is particularly important in the distributed construction, where it reduces the communication cost.

We discover that for cases where we have a cluster node of GPUs, the communication cost is negligible, and the distributed construction is instead dominated by the calculation of peer ranks. We study the effect of different algorithm parameters on performance. We optimize the local construction method and achieve performance improvements up to $2.45 \times$ for 100 million particles.

\end{abstract}
